\documentclass[11pt,a4paper]{article}

% --- Pacchetti ---
\usepackage[utf8]{inputenc}
\usepackage[T1]{fontenc}
\usepackage{lmodern}
\usepackage[italian]{babel}
\usepackage[margin=1in]{geometry}
\usepackage{xcolor}
\usepackage{listings}
\usepackage{hyperref}
\usepackage{enumitem}

% --- Colori ---
\definecolor{primary}{RGB}{38, 66, 139}
\definecolor{codebg}{RGB}{245, 245, 245}

\hypersetup{colorlinks=true, linkcolor=primary, urlcolor=primary}

% --- Configurazione Codice ---
\lstset{
    backgroundcolor=\color{codebg},
    basicstyle=\ttfamily\small,
    breaklines=true,
    frame=single,
    keywordstyle=\color{blue},
    commentstyle=\color{green!50!black},
    stringstyle=\color{red!70!black},
    language=Python
}

\title{\textbf{\color{primary} Esercitazione: Creazione di un Bot Webex}}
\author{Programmazione e Tecnologie Web - Lezione 3}
\date{A.A. 2025/2026}

\begin{document}

\maketitle

\section{Obiettivo}
Mettere in pratica i concetti di \textbf{Web API RESTful} e \textbf{Bearer Token} creando un'applicazione Python capace di interagire con la piattaforma Cisco Webex.

\section{Prerequisiti}
\begin{itemize}
    \item Python 3.x e libreria \texttt{requests} (\texttt{pip install requests}).
    \item Account su \href{https://developer.webex.com/}{developer.webex.com}.
\end{itemize}

\section{Esercizio Preliminare: Esplorazione con Postman / VS Code}
Prima di scrivere codice, è fondamentale testare le API manualmente utilizzando dei \textit{REST Clients}. Questo permette di visualizzare immediatamente gli header e la struttura JSON di risposta.

\subsection{Utilizzo di Postman}
\begin{enumerate}
    \item Scaricare e aprire \textbf{Postman}.
    \item Creare una nuova richiesta di tipo \texttt{GET} verso: \url{https://jsonplaceholder.typicode.com/posts/1}.
    \item Cliccare su \textbf{Send} e osservare il corpo della risposta (JSON) e lo Status Code (\texttt{200 OK}).
    \item Provare a inviare una richiesta verso un indirizzo errato (es. aggiungere \texttt{/xyz} alla fine) e osservare il \texttt{404 Not Found}.
\end{enumerate}

\subsection{Utilizzo di VS Code (Thunder Client / REST Client)}
Se preferite rimanere all'interno dell'editor, potete installare un'estensione dedicata:
\begin{itemize}
    \item \textbf{Thunder Client}: Molto simile a Postman, permette di gestire collezioni di richieste tramite una GUI integrata in VS Code.
    \item \textbf{REST Client}: Permette di scrivere le richieste in un semplice file con estensione \texttt{.http} (es. \texttt{test.http}).
\end{itemize}

\textbf{Task}: Effettuare una richiesta \texttt{GET} verso \texttt{/v1/people/me} delle API Webex (endpoint: \url{https://webexapis.com/v1/people/me}). 
\begin{itemize}
    \item Dovrete inserire il vostro \texttt{ACCESS\_TOKEN} nel tab \textbf{Authorization} selezionando il tipo \textbf{Bearer Token}.
    \item Verificate che il server restituisca correttamente i vostri dati personali (nome, cognome, email).
\end{itemize}

\section{Step 1: Ottenere il Personal Access Token}
Per autenticarsi presso le API di Webex, è necessario un token. 
\begin{enumerate}
    \item Effettuare il login su \url{https://developer.webex.com/}.
    \item Navigare in \textbf{Documentation} $\rightarrow$ \textbf{APIs} $\rightarrow$ \textbf{Messages} $\rightarrow$ \textbf{Create a Message}.
    \item Nella sezione "Try it" sulla destra, assicurarsi che "Personal Access Token" sia selezionato e copiare la stringa alfanumerica (\textit{Bearer Token}).
\end{enumerate}

\section{Step 2: Implementazione dello Script}
Completare il seguente schema nel file \texttt{bot.py}:

\begin{lstlisting}
import requests

# CONFIGURAZIONE
ACCESS_TOKEN = "IL_TUO_TOKEN_QUI"
DESTINATARIO = "email@esempio.com" 

def invia_messaggio(testo):
    url = "https://webexapis.com/v1/messages"
    
    headers = {
        "Authorization": f"Bearer {ACCESS_TOKEN}",
        "Content-Type": "application/json"
    }
    
    payload = {
        "toPersonEmail": DESTINATARIO,
        "text": testo
    }
    
    response = requests.post(url, headers=headers, json=payload)
    
    if response.status_code == 200:
        print("Successo!")
    else:
        print(f"Errore: {response.status_code}")

invia_messaggio("Hello from Python!")
\end{lstlisting}

\section{Esercizi Suggeriti}
\begin{enumerate}
    \item \textbf{Input Dinamico}: Modificare lo script affinché il testo del messaggio venga inserito dall'utente tramite terminale.
    \item \textbf{Gestione Errori}: Utilizzare \texttt{response.raise\_for\_status()} e un blocco \texttt{try/except} per gestire eventuali problemi di rete o token scaduti.
    \item \textbf{GET Request}: Creare una funzione \texttt{leggi\_messaggi()} che interroghi l'endpoint \texttt{GET /v1/messages} per stampare gli ultimi messaggi della chat.
\end{enumerate}

\end{document}
